% Source: dropout_explanation.tex
\begin{notechapter}
% ---- chapter-local definitions ----
\lstset{
    language=Python,
    basicstyle=\ttfamily\small,
    keywordstyle=\color{blue},
    commentstyle=\color{gray},
    stringstyle=\color{red},
    showstringspaces=false,
    breaklines=true,
    frame=single,
    numbers=left,
    numberstyle=\tiny\color{gray}
}
% -----------------------------------
\chapter{Dropout: From Theory to Implementation}\label{ch:dropout-inference}

\section{Introduction}

Dropout is a regularization technique for neural networks that randomly removes neurons during training. This document explains the mathematical foundations of dropout, particularly focusing on how inference is performed and why the layer-by-layer approximation works.

\section{Dropout During Training}

\begin{definition}[Dropout Layer]
Let $\mathbf{X} \in \mathbb{R}^{n \times c}$ be a mini-batch of activations with $n$ examples and $c$ features. A dropout layer with keep probability $p \in [0, 1]$ applies:
\begin{equation}
\text{Dropout}(\mathbf{X}) = \mathbf{M} \odot \mathbf{X}
\end{equation}
where $\mathbf{M} \in \{0, 1\}^{n \times c}$ is a binary mask with $M_{ij} \sim \text{Bernoulli}(p)$ sampled independently.
\end{definition}

\begin{remark}
In this chapter \( p \) is the \emph{keep} probability; \( 1 - p \) is the \emph{drop probability}, the fraction of neurons that are zeroed out. Beware that deep-learning libraries (e.g.\ \texttt{torch.nn.Dropout(p)}) and the previous chapter use \( p \) for the \emph{drop} probability.
\end{remark}

\subsection{Why Dropout Works}

During training, dropout:
\begin{itemize}
    \item Forces the network to learn redundant representations
    \item Prevents co-adaptation of neurons (where neurons rely too heavily on specific other neurons)
    \item Can be viewed as training an ensemble of $2^n$ subnetworks with shared weights
    \item Adds noise to intermediate representations, similar to data augmentation
\end{itemize}

\section{The Challenge of Inference}

After training, we want deterministic predictions. However, the output of a network with dropout is a random variable due to the random masks $\mathbf{M}$.

\subsection{Problem Statement}

Consider a network $f(\mathbf{x}; \mathbf{M})$ where $\mathbf{M} = \{M_1, M_2, \ldots, M_L\}$ represents all dropout masks across $L$ layers. For a given input $\mathbf{x}$, the output is stochastic:
\begin{equation}
\mathbf{y} = f(\mathbf{x}; \mathbf{M}), \quad \mathbf{M} \sim p(\mathbf{M})
\end{equation}

We need a deterministic function for inference.

\section{Monte Carlo Dropout}

\subsection{Full Expectation Approach}

The ideal deterministic output would be:
\begin{equation}
\mathbb{E}_{p(\mathbf{M})}[f(\mathbf{x}; \mathbf{M})] = \int f(\mathbf{x}; \mathbf{M}) \, p(\mathbf{M}) \, d\mathbf{M}
\end{equation}

Since we cannot compute this integral analytically, we can approximate it via Monte Carlo sampling:
\begin{equation}
\mathbb{E}_{p(\mathbf{M})}[f(\mathbf{x}; \mathbf{M})] \approx \frac{1}{K} \sum_{k=1}^{K} f(\mathbf{x}; \mathbf{M}^{(k)}), \quad \mathbf{M}^{(k)} \sim p(\mathbf{M})
\end{equation}

\begin{remark}
This approach requires $K$ forward passes through the network, which can be computationally expensive. However, it provides uncertainty estimates and is sometimes used in Bayesian deep learning.
\end{remark}

\section{Layer-by-Layer Approximation}

\subsection{The Key Insight}

Instead of computing the expectation over all masks simultaneously, we approximate by replacing each dropout layer's mask with its expectation \emph{independently}.

\subsection{Mathematical Derivation}

For a single dropout layer with input $\mathbf{X}$:
\begin{align}
\mathbb{E}_{p(M)}[\text{Dropout}(\mathbf{X})] &= \mathbb{E}_{p(M)}[\mathbf{M} \odot \mathbf{X}] \\
&= \mathbb{E}[\mathbf{M}] \odot \mathbf{X} \quad \text{(since $\mathbf{X}$ is deterministic)}
\end{align}

Computing element-wise:
\begin{align}
\mathbb{E}[M_{ij} \cdot X_{ij}] &= X_{ij} \cdot \mathbb{E}[M_{ij}] \\
&= X_{ij} \cdot p \quad \text{(since $M_{ij} \sim \text{Bernoulli}(p)$)}
\end{align}

Therefore:
\begin{equation}
\boxed{\mathbb{E}_{p(M)}[\text{Dropout}(\mathbf{X})] = p\mathbf{X}}
\end{equation}

\subsection{Why ``Layer-by-Layer''?}

For a network with multiple dropout layers:

\begin{itemize}
    \item \textbf{Exact expectation:} 
    \begin{equation}
    \mathbb{E}_{p(M_1, M_2, \ldots, M_L)}[f(\mathbf{x}; M_1, M_2, \ldots, M_L)]
    \end{equation}
    This requires integrating over all possible combinations of masks across all layers.
    
    \item \textbf{Layer-by-layer approximation:}
    \begin{equation}
    f(\mathbf{x}; \mathbb{E}[M_1], \mathbb{E}[M_2], \ldots, \mathbb{E}[M_L]) = f(\mathbf{x}; p_1\mathbf{I}, p_2\mathbf{I}, \ldots, p_L\mathbf{I})
    \end{equation}
    We replace each mask independently with its expectation.
\end{itemize}

\begin{remark}
This is an approximation because:
\begin{equation}
\mathbb{E}[g(h(\mathbf{X}))] \neq g(\mathbb{E}[h(\mathbf{X})]) \quad \text{in general}
\end{equation}
The expectation does not commute with nonlinear compositions. However, this approximation works well in practice and is computationally efficient.
\end{remark}

\section{Inverted Dropout}

\subsection{Motivation}

The formulation above requires scaling by $p$ during inference. An alternative is \emph{inverted dropout}, which performs the scaling during training instead.

\subsection{Standard vs. Inverted Dropout}

\begin{center}
\begin{tabular}{|l|c|c|}
\hline
& \textbf{Training} & \textbf{Inference} \\
\hline
\textbf{Standard Dropout} & $\mathbf{M} \odot \mathbf{X}$ & $p\mathbf{X}$ \\
\hline
\textbf{Inverted Dropout} & $\frac{1}{p}(\mathbf{M} \odot \mathbf{X})$ & $\mathbf{X}$ \\
\hline
\end{tabular}
\end{center}

\subsection{Why They're Equivalent}

For inverted dropout during training:
\begin{align}
\mathbb{E}\left[\frac{1}{p}(\mathbf{M} \odot \mathbf{X})\right] &= \frac{1}{p} \mathbb{E}[\mathbf{M} \odot \mathbf{X}] \\
&= \frac{1}{p} \cdot p\mathbf{X} \\
&= \mathbf{X}
\end{align}

During inference, we simply use $\mathbf{X}$ (no scaling needed). The expected values match in both training and inference!

\begin{proposition}
Inverted dropout ensures that the expected output of a layer is the same during training and inference, eliminating the train-test mismatch.
\end{proposition}

\section{Implementation in PyTorch}

\subsection{The Role of \texttt{model.train()} and \texttt{model.eval()}}

PyTorch implements inverted dropout. The behavior changes based on the model's mode:

\begin{lstlisting}[language=Python]
import torch
import torch.nn as nn

# Define a simple model with dropout
model = nn.Sequential(
    nn.Linear(10, 20),
    nn.ReLU(),
        nn.Dropout(p=0.5),  # PyTorch's p = DROP probability (here 50% kept)
    nn.Linear(20, 1)
)

x = torch.randn(16, 10)  # Batch of 16 samples

# Training mode: dropout is active
model.train()
y_train = model(x)  # Output is stochastic

# Evaluation mode: dropout is disabled
model.eval()
y_eval = model(x)  # Output is deterministic

# The outputs will be different!
print("Training output:", y_train[0])
print("Eval output:", y_eval[0])
\end{lstlisting}

\subsection{What Happens Internally}

When \texttt{model.train()} is called:
\begin{itemize}
    \item Sets \texttt{self.training = True} on all layers
        \item Dropout layer applies: $\text{output} = \frac{1}{p}(\mathbf{M} \odot \mathbf{X})$ with $p$ the keep probability of this chapter; in terms of PyTorch's argument $p_{\text{drop}}$, the scaling factor is $\frac{1}{1 - p_{\text{drop}}}$
\end{itemize}

When \texttt{model.eval()} is called:
\begin{itemize}
    \item Sets \texttt{self.training = False} on all layers  
    \item Dropout layer applies: $\text{output} = \mathbf{X}$ (identity operation)
\end{itemize}

\subsection{Monte Carlo Dropout in PyTorch}

If you want to use Monte Carlo dropout for uncertainty estimation:

\begin{lstlisting}[language=Python]
# Keep dropout active during inference
model.train()

# Multiple forward passes
K = 10
x = torch.randn(1, 10)  # Single input
predictions = []

for _ in range(K):
    y = model(x)
    predictions.append(y)

# Average predictions
mean_pred = torch.stack(predictions).mean(dim=0)
std_pred = torch.stack(predictions).std(dim=0)

print(f"Mean prediction: {mean_pred}")
print(f"Uncertainty (std): {std_pred}")
\end{lstlisting}

\section{Practical Considerations}

\subsection{When to Use Dropout}

Dropout is most commonly used for:
\begin{itemize}
    \item Fully-connected (dense) layers
    \item Attention layers in transformers
    \item Less common for convolutional layers (use spatial dropout instead)
\end{itemize}

\subsection{Typical Keep Probabilities}

\begin{itemize}
    \item $p = 0.8$ or $p = 0.9$ for fully-connected layers
    \item $p = 0.1$ to $p = 0.3$ for attention layers
    \item Often no dropout for the first and last layers
\end{itemize}

\subsection{Alternatives for CNNs}

For convolutional layers, standard dropout is too unstructured. Better alternatives:

\begin{itemize}
    \item \textbf{Spatial Dropout:} Drop entire channels instead of individual elements
    \item \textbf{Cutout:} Drop spatial patches (rectangular regions)
    \item \textbf{DropConnect:} Drop weights instead of activations
\end{itemize}

\section{Summary}

\begin{enumerate}
    \item Dropout randomly zeros out neurons during training with probability $1-p$
    \item During inference, we need a deterministic operation
    \item The layer-by-layer approximation replaces each dropout mask with its expectation independently: $\mathbb{E}[\mathbf{M}] = p\mathbf{I}$
    \item PyTorch uses inverted dropout: scale by $1/p$ during training, do nothing during inference
    \item \texttt{model.train()} enables dropout; \texttt{model.eval()} disables it
    \item Monte Carlo dropout keeps dropout active during inference for uncertainty estimation
\end{enumerate}

\section*{Key Equation}

The fundamental result for the layer-by-layer approximation:
\begin{equation}
\boxed{\mathbb{E}_{p(M)}[\text{Dropout}(\mathbf{X})] = p\mathbf{X}}
\end{equation}

This means we can replace the stochastic dropout operation during inference with a simple deterministic scaling by $p$ (or in inverted dropout, no operation at all).
\end{notechapter}
